\section{Discussion and Bibliography}
\label{sec:1.6}

One of the earliest probabilistic experiments resembling modern Monte Carlo methods is Buffon's
needle experiment, proposed and solved by Georges-Louis Leclerc, Comte de Buffon, in the eighteenth
century~\cite{51}. Later work~\cite{134} observed that repeated trials of this experiment yield a
stochastic procedure for approximating $\pi$. Historically, a major limitation of Monte Carlo
methods was the cost of carrying out random experiments and processing their results. Physical
randomization procedures were difficult to perform repeatedly, and the numerical calculations
required to analyze their outcomes were often laborious. As an early attempt to streamline this
process, Enrico Fermi studied neutron diffusion using a mechanical device known as the
\textit{Fermiac}, which simulated random neutron trajectories through matter.

The situation changed fundamentally with the advent of the digital computer. Among the first to
realize its potential for stochastic simulation were John von Neumann and Stanislaw Ulam, who were
then working (along with Fermi) for the Manhattan Project at Los Alamos National Laboratory; see
\cite{67}. The development of Monte Carlo simulation required reliable methods for generating random
numbers computationally. Early work by von Neumann introduced simple pseudo-random number
generators~\cite{232}, illustrating how randomness could be produced algorithmically rather than
through physical processes. In 1947 Ulam and von Neumann proposed using computer simulations to
study neutron diffusion in fissionable materials. The 1949 paper by Nicholas Metropolis and
Ulam~\cite{163} introduced the first unclassified, formal description of the Monte Carlo method.

At that time ---and for several decades thereafter--- programming was often regarded as routine
technical work and was frequently performed by women. Kl\'{a}ra D\'{a}n von Neumann, the wife of
John von Neumann, is credited with writing the code for one of the first computer implementations
of a Monte Carlo simulation. Interestingly, John and Kl\'{a}ra von Neumann first met in a casino in
Monte Carlo. However, the name ``Monte Carlo'' is generally attributed to Ulam, who reportedly drew
inspiration from his interest in games of chance and from an uncle who frequently borrowed money to
gamble in Monte Carlo (unpublished remarks by Ulam, 1983).

Important methodological advances during this early period included the development of variance
reduction techniques and importance sampling, which significantly improved the efficiency of Monte
Carlo estimators. A major breakthrough occurred in 1953 when Metropolis and collaborators calculated
the equation of state via modified Monte Carlo integration over configuration space~\cite{162}. This
work established the basic idea of constructing a Markov chain whose stationary distribution
coincides with the target distribution of interest. The method was later generalized by Wilfred
Keith Hastings in 1970, resulting in what is now known as the Metropolis Hastings
algorithm~\cite{108}.

Despite these advances, Monte Carlo methods were not widely used in mainstream statistics until the
1980s and 1990s. A major turning point came with the introduction of the Gibbs sampling method by
Stuart Geman and Donald Geman~\cite{82}. Soon afterward, work by Alan Gelfand and Adrian
Smith~\cite{80} demonstrated how Markov chain Monte Carlo methods could be used to perform Bayesian
inference in complex statistical models, thereby triggering a rapid expansion of Monte Carlo
methodology within statistics.

Since then, the field has continued to develop in several directions. Sequential Monte Carlo
methods, or particle filters, were introduced for inference in dynamical systems and state-space
models~\cite{54,49}. In parallel, Quasi-Monte Carlo methods were developed to improve numerical
integration by replacing random samples with carefully constructed low-discrepancy
sequences~\cite{174}, and Monte Carlo algorithms based on Hamiltonian and Langevin dynamics grew in
importance~\cite{60,171,20,203}. More recently, the core ideas underlying Monte Carlo methods have
found applications in generative modeling and modern machine learning, where stochastic simulation
and probabilistic inference play a central role.
